\documentclass{article}
\usepackage{mathtools} 
\usepackage{fontspec}
\usepackage[UTF8]{ctex}
\usepackage{amsthm}
\usepackage{mdframed}
\usepackage{xcolor}
\usepackage{amssymb}
\usepackage{amsmath}
\usepackage{tikz}
\usepackage{pgfplots}
\usepgfplotslibrary{polar}
\usetikzlibrary{3d} % 允许3D坐标
% \pgfplotsset{compat=1.18}


% 定义新的带灰色背景的说明环境 zremark
\newmdtheoremenv[
  backgroundcolor=gray!10,
  % 边框与背景一致，边框线会消失
  linecolor=gray!10
]{zremark}{说明}

% 通用矩阵命令: \flexmatrix{矩阵名}{元素符号}{行数}{列数}
\newcommand{\flexmatrix}[4]{
  \[
  #1 = \begin{pmatrix}
    #2_{11}     & #2_{12}     & \cdots & #2_{1#4}   \\
    #2_{21}     & #2_{22}     & \cdots & #2_{2#4}   \\
    \vdots      & \vdots      & \ddots & \vdots     \\
    #2_{#31}    & #2_{#32}    & \cdots & #2_{#3#4}
  \end{pmatrix}
  \]
}

% 简化版命令（默认矩阵名为A，元素符号为a）: \quickmatrix{行数}{列数}
\newcommand{\quickmatrix}[2]{\flexmatrix{A}{a}{#1}{#2}}

\begin{document}
\title{习题 10.4}
\author{张志聪}
\maketitle

\section*{2}

提示：极坐标改为直角坐标，利用参数方程求旋转曲面的面积。

\begin{equation*}
  \begin{cases}
    x & = r(\theta) cos \theta \\
    y & = r(\theta) sin \theta
  \end{cases}
\end{equation*}
于是
\begin{align*}
  x' & = r'cos \theta - r sin \theta \\
  y' & = r'sin \theta + r cos \theta
\end{align*}
由旋转曲面的参数方程可知
\begin{align*}
  S & = 2\pi \int_{0}^{\pi} y(\theta) \sqrt{x'^2(\theta) + y'^2(\theta)} d\theta                                                                                                      \\
    & = 2\pi \int_{0}^{\pi} r sin \theta \sqrt{[r' cos \theta - r sin \theta]^2 + [r' sin \theta + r cos \theta]} d\theta                                                             \\
    & = 2\pi \int_{0}^{\pi} r sin \theta \sqrt{r'^2 cos^2 \theta + r^2 sin^2 \theta - 2 r'r sin\theta cos \theta + r'^2 sin^2 \theta + r^2 cos^2 \theta + 2 r'r sin\theta cos \theta} \\
    & = 2\pi \int_{0}^{\pi} r sin \theta \sqrt{r^2 + r'^2} d \theta
\end{align*}

\section*{3}

\begin{itemize}
  \item (2)

        \begin{tikzpicture}[scale=3]
          \def\a{1}
          \draw[->] (-1.5,0)--(1.5,0);
          \draw[->] (0,-1.5)--(0,1.5);
          \draw[domain=-45:45,smooth,variable=\t,blue,thick]
          plot({sqrt(2)*\a*sqrt(cos(2*\t))*cos(\t)},
          {sqrt(2)*\a*sqrt(cos(2*\t))*sin(\t)});
          \draw[domain=-45:45,smooth,variable=\t,blue,thick]
          plot({-sqrt(2)*\a*sqrt(cos(2*\t))*cos(\t)},
          {-sqrt(2)*\a*sqrt(cos(2*\t))*sin(\t)});
          \def\theta{30}
          \def\r{sqrt(2)*\a*sqrt(cos(2*\theta))}
          \fill[red] ({\r*cos(\theta)},{\r*sin(\theta)}) circle(0.8pt);
          \fill[red] ({-\r*cos(\theta)},{-\r*sin(\theta)}) circle(0.8pt);
          \draw[dashed,gray] (0,0)--({1.5*cos(\theta)},{1.5*sin(\theta)});
          \draw[dashed,gray] (0,0)--({-1.5*cos(\theta)},{-1.5*sin(\theta)});
        \end{tikzpicture}

        通过画图流程可知，由对称性，我们有
        \begin{align*}
          S = 2 \cdot 2 \pi \int_{0}^{\frac{\pi}{4}} r sin \theta \sqrt{r^2 + r'^2} d \theta
        \end{align*}
        计算略。

\end{itemize}

\end{document}